% Contents/inventario_e_carga.tex
\midsectiontitle{Inventário e Carga}

\section*{O Fardo do Aventureiro}

No universo de \textit{Overgrown}, cada objeto tem sua propria unidade de peso. A capacidade de carga básica de um indivíduo é de \textbf{10 unidades}, representando o limite de volume e peso que pode ser manejado sem comprometer a mobilidade.

Este limite é ampliado pela força física do aventureiro: para cada ponto investido em \link{att:might}{Might (Poder)}, o aventureiro expande sua capacidade de carga em \textbf{1 unidade}. Embora o Pináculo não imponha um limite absoluto para o que se pode carregar, o excesso de peso além da capacidade total resultará em penalidades severas de movimento e ações. Pra cada 1 unidade de Peso acima da carga do aventureiro, este perde 5 de movimento.

\begin{rpgboxtips}{Bonificações Efêmeras}
Bônus temporários de atributos ou estados de fúria momentânea não expandem sua capacidade de carga máxima. Apenas melhorias permanentes provenientes da \textbf{Árvore de Talentos} ou artefatos  alteram o limite de transporte do personagem.
\end{rpgboxtips}

\section*{Armamentos e Projéteis}

O manejo de recursos é vital para a sobrevivência. Armas que dependem de munição seguem regras distintas baseadas em sua natureza.

\subsection*{A Ciência da Pólvora}
Armas de fogo, como pistolas, rifles e fuzis são maravilhas da engenharia rúnica. Elas não são meramente mecânicas, mas forjadas com gravuras que estabilizam a combustão da \textbf{Munição de Pólvora}. 
\begin{itemize}
    \item \textbf{Consumo:} Cada unidade de munição de pólvora é preparada para sustentar o aventureiro por uma \textbf{Cena} inteira de combate.
    \item \textbf{Carga:} Cada suprimento de cena pesa \textbf{1 unidade}.
    \item \textbf{Natureza:} Embora sejam armas mágicas por fabricação, elas são focadas na entrega de dano físico bruto e, por regra geral, não servem como canalizadores para magias ativas.
\end{itemize}

\subsection*{A Arte do Arco e Flecha}
Diferente das armas de fogo, os arcos são extensões da alma do arqueiro e excelentes condutores de energia primordial. 
\begin{itemize}
    \item \textbf{Consumo:} Flechas básicas não são contabilizadas por combate, mas sim por \textbf{Aventura}, tendo que ser repostas em períodos de descanso ou em mercadores.
    \item \textbf{Canalização:} Arcos podem ser usados como canalizadores mágicos, permitindo disparar feitiços como se fossem flechas. Nesta forma, a magia herda as propriedades do arco, como o \link{stats:crit}{Nível de Ameaça (Crítico)} e bônus de distância.
\end{itemize}

\begin{center}
\renewcommand{\arraystretch}{1.2}
\begin{tabular}{lc}
\toprule
\rowcolor{gray!15} \textbf{\ringm{Tipo de Projétil}} & \textbf{\ringm{Peso de Carga}} \\ 
\midrule
Flechas de Ferro ou Aço & 1 \\
Flechas de Aetherium (Vastidrum) & 2 \\
Flechas Rúnicas (Engravadas) & 3 \\
\bottomrule
\end{tabular}
\end{center}

\section*{Proteção e Utilitários}

As armaduras são classificadas por sua densidade e material, afetando diretamente o fardo do usuário. Armaduras especiais podem ter pesos alterados conforme suas especificações técnicas.

\begin{description}
    \item[\ringm{Pano (Peso 1)}:] Vestes leves geralmente usada por magos.
    \item[\ringm{Couro (Peso 2)}:] Proteção flexível e leve.
    \item[\ringm{Malha (Peso 4)}:] O equilíbrio entre defesa e peso.
    \item[\ringm{Metal (Peso 6)}:] Proteção pesada para a linha de frente.
    \item[\ringm{Placas (Peso 8)}:] A fortaleza intransponível do guerreiro.
\end{description}

Alguns itens de consumo imediato, como poções de cura, são leves o suficiente para não ocuparem carga individualmente. No entanto, ferramentas complexas como \textbf{Kits Médicos} ocupam \textbf{1 unidade} de carga devido à sua composição de suprimentos.

\begin{rpgboxtips}{Conselho aos Mestres}
Utilize a lista de \link{data:itens}{itens preexistentes} como base de equilíbrio. Ao criar novos artefatos, compare sua utilidade e material com as categorias acima para atribuir um peso justo.
\end{rpgboxtips}